% The defect this chapter prevents, drawn to scale: delaying for a period
% against a timer that owns the period. The tick pitch is 1.65 and the wrong
% way's pitch is 1.80, so the drift in the top row is the figure's own subject
% rather than an artefact of the drawing.
\begin{tikzpicture}[font=\sffamily\scriptsize, x=1cm, y=1cm]

% ================================================================ the wrong way
\node[signame, anchor=east] at (-0.3,2.35) {delay for 1 ms};
\foreach \i in {0,1,2,3,4,5}{ \draw[tick] (\i*1.65,1.55) -- (\i*1.65,3.15); }

\foreach \x in {0, 1.80, 3.60, 5.40}{
  \node[taskrun, minimum width=6mm, anchor=south west] at (\x,2.0) {work};
  \draw[tspan] (\x+0.62,2.18) -- node[tlabel]{1 ms} (\x+1.80,2.18);
}
\node[taskrun, minimum width=6mm, anchor=south west] at (7.20,2.0) {work};

\foreach \i in {1,2,3,4}{
  \draw[faultedge, <->] (\i*1.65,1.72) -- (\i*1.65 + \i*0.15,1.72);
}
\node[tlabel, anchor=west, text=red!60!black] at (8.1,1.72)
  {the error accumulates without limit};

% ================================================================ the right way
\node[signame, anchor=east] at (-0.3,0.35) {timer owns the period};
\foreach \i in {0,1,2,3,4,5}{
  \draw[tick] (\i*1.65,-0.62) -- (\i*1.65,1.15);
  \draw[release] (\i*1.65,-0.52) -- (\i*1.65,0.02);
}
\foreach \i/\d in {0/0.00, 1/0.07, 2/0.02, 3/0.16, 4/0.04, 5/0.09}{
  \node[taskrun, minimum width=6mm, anchor=south west] at (\i*1.65 + \d,0.0) {work};
}
\foreach \i in {0,1,2,3,4}{
  \draw[tspan] (\i*1.65,-0.3) -- node[tlabel]{1 ms, exactly} (\i*1.65+1.65,-0.3);
}
\node[tlabel, anchor=west] at (8.4,0.25) {release jitter only};

% ================================================================ overrun
\node[signame, anchor=east] at (-0.3,-1.75) {work overruns};
\foreach \i in {0,1,2,3,4,5}{ \draw[tick] (\i*1.65,-2.4) -- (\i*1.65,-1.3); }
\node[taskrun, minimum width=6mm, anchor=south west] at (0,-2.1) {work};
\node[taskrun, minimum width=6mm, anchor=south west] at (1.65,-2.1) {work};
\node[taskrun, minimum width=21mm, anchor=south west, fill=hwcol] at (3.30,-2.1) {a long one};
\node[taskrun, minimum width=6mm, anchor=south west] at (6.60,-2.1) {work};
\draw[faultedge] (4.95,-1.3) -- (4.95,-1.78);
\node[tlabel, anchor=south, text=red!60!black] at (4.95,-1.28) {tick missed};
\node[tlabel, anchor=west] at (8.4,-1.85) {counted, reported, not hidden};

% ================================================================ the reading
\node[umlnote, text width=52mm, anchor=north west] at (-1.0,-3.1)
  {The top two rows look identical on an instrument that samples every ten
   microseconds, and they are not the same machine. The top row runs at one
   millisecond plus the work, forever; after an hour it has lost a number of
   periods equal to the total work time divided by the period. The middle row
   runs at one millisecond and moves only the start of the work, which is
   release jitter and is what the histogram measures.};
\node[umlnote, text width=44mm, anchor=north east] at (11.2,-3.1)
  {The bottom row is the case the budget exists for. A handler still working
   when the next tick arrives has missed a deadline, and the only correct
   response is to count it and carry on. A node that silently runs at half rate
   is worse than one that reports.};
\end{tikzpicture}
